\documentclass[a5paper]{ltjsarticle}
\begin{document}
\section{図と表}
朝から細かい雨が降っていた。傘を持たずに家を出たのは失敗だったと、駅までの坂道を歩きながら何度も思った。商店街の軒先を伝って歩けば少しは濡れずに済むが、それでも肩のあたりはすぐに冷たくなった。

\begin{figure}[t]
\centering
\rule{6\zw}{3\zw}
\caption{上に置いた図}
\end{figure}

電車は思ったより空いていた。窓際の席に座ると、外の景色が雨粒で滲んで見えた。友人は鞄から文庫本を取り出し、しおりを挟んだ頁を開いた。

\begin{table}[b]
\centering
\caption{下に置いた表}
\begin{tabular}{|c|c|c|}\hline
項目 & 数 & 備考\\\hline
りんご & 三 & 赤い\\\hline
みかん & 五 & 甘い\\\hline
\end{tabular}
\end{table}

例の場所というのが何を指すのか、私にはすぐにわかった。学生の頃によく通った川沿いの喫茶店である。店主は無口な人で、注文を聞く以外にはほとんど口をきかなかった。

縦組みの表を本文中に置く：
\begin{center}
\begin{tabular}{|l|l|}\hline
一行目 & 縦組みの表\\\hline
二行目 & 文字が縦に並ぶ\\\hline
\end{tabular}
\end{center}
横組みのまま置いた表：
\begin{center}
\hbox{\yoko\begin{tabular}{|l|l|}\hline
A & 横組みの表\\\hline
B & 文字は横\\\hline
\end{tabular}}
\end{center}
表のあとの本文。
\end{document}
